% Complete proofs for strict single-peaked majority and the median conclusions.
\section{Single-peaked majority: proofs}
\label{sec:annex-p17-single-peaked-majority}

The source's interval exercise is Rubinstein, Lecture 9, Problem 2
\cite[p.~128]{Rubinstein2015LectureNotes}. The strict input and weak output
interfaces come from Chapter~\ref{sec:p17-social-choice}. The common-axis
definitions used in these proofs are in
Chapter~\ref{sec:p17-single-peaked-majority}.

\begin{annexitem}{annex:p17.single-peaked-structure}{thm:p17-single-peaked-structure}{Peaks and the absence of a middle-ranked valley}
\begin{proof}
If \(x<p_i\), take \(b=p_i,c=x\) in the single-peaked condition;
it gives \(p_i\succ_i x\). If \(x>p_i\), take \(b=p_i,c=x\)
in its other branch, obtaining the same conclusion. These are all
possibilities for \(x\ne p_i\). If a different point also satisfied
the peak condition, each of the two peaks would be strictly preferred
to the other, contrary to asymmetry of a strict ordering. Thus the peak
is unique and is the greatest preferred alternative.

Fix \(a<b<c\). If \(b\leq p_i\), the first branch of the
condition gives \(b\succ_i a\). If \(p_i\leq b\), the second
branch gives \(b\succ_i c\). A linear axis guarantees at least
one of these inequalities between \(b,p_i\), including both when
\(b=p_i\). Hence \(b\) cannot be below both endpoints. Of the six
strict orderings of three distinct alternatives, exactly
\(a\succ_i c\succ_i b\) and \(c\succ_i a\succ_i b\) put \(b\)
last. Excluding them leaves exactly the four listed forms.

For the converse, assume \(p\) is greatest under the strict ordering
and every axis-ordered triple has the stated no-valley property.
If \(p\geq b>c\) and \(b=p\), then \(b\succ c\) by maximality.
If instead \(c<b<p\), failure of \(b\succ c\) would imply
\(c\succ b\), because the ordering compares distinct points.
Also \(p\succ b\). This makes the middle member \(b\) of
\(c<b<p\) worse than both endpoints, a contradiction.
Thus \(b\succ c\) throughout this first branch.
If \(c>b\geq p\), the case \(b=p\) is again maximality;
otherwise \(p<b<c\), and failure of \(b\succ c\) would put
\(b\) below both \(p,c\), the same forbidden pattern. This proves
both branches of single-peakedness.

A strict ordering on a finite nonempty set has a greatest element:
start from any element and compare it successively with each other
element, retaining the better one. At each step the retained element
is above all previously discarded ones by transitivity. The final
element is above every other one. This establishes the finite-set
qualification in the converse without assuming a greatest element on
an arbitrary infinite set.
\end{proof}
\end{annexitem}

\begin{annexitem}{annex:p17.single-peaked-majority}{thm:p17-single-peaked-majority}{Excluding each three-cycle and proving full transitivity}
\begin{proof}
For each distinct pair, strict individual preferences give
\(V_P(x,y)+V_P(y,x)=2k+1\). These integer counts cannot be equal.
Exactly one is at least \(k+1\), so exactly one strict majority
direction holds. In particular, the relation is asymmetric and has
no self-comparison. Its weak version is its reflexive closure.

Consider any three distinct alternatives, and name them \(a<b<c\)
according to their positions on the common axis. The preceding lemma's
four-form list implies two inclusions of voter sets:
\[
 \{i:a\succ_i b\}\subseteq\{i:a\succ_i c\},\qquad
 \{i:c\succ_i b\}\subseteq\{i:c\succ_i a\}.
\]
Indeed, the only allowed ordering that puts \(a\) above \(b\)
is \(a\succ_i b\succ_i c\); the only one putting \(c\) above
\(b\) is \(c\succ_i b\succ_i a\).
Counting the first inclusion proves
\(a\succ^M_P b\Rightarrow a\succ^M_P c\).
Counting the second proves
\(c\succ^M_P b\Rightarrow c\succ^M_P a\).

A directed majority cycle on these three alternatives has one of two
orientations. If its comparison of \(a,b\) is \(a\succ^M_P b\),
the cycle would have \(b\succ^M_P c\) and \(c\succ^M_P a\).
The first implication just proved instead gives \(a\succ^M_P c\),
contradicting asymmetry. If its comparison is \(b\succ^M_P a\),
the cycle would have \(a\succ^M_P c\) and \(c\succ^M_P b\).
The second implication gives \(c\succ^M_P a\), the opposite
contradiction. Thus neither three-cycle is possible.

Now suppose \(x\succ^M_P y\) and \(y\succ^M_P z\).
The three points must be distinct: each strict comparison excludes
equality of its arguments, and \(x=z\) would contradict asymmetry.
If \(x\succ^M_P z\) were false, the strict comparison of that
pair would give \(z\succ^M_P x\), producing a prohibited
three-cycle. Therefore \(x\succ^M_P z\). This is transitivity
of the entire majority relation on \(X\), including an infinite \(X\).
It does not infer transitivity from the existence of a greatest
alternative or restrict attention to comparisons with such an alternative.

The reflexive closure is complete. To check weak transitivity, suppose
\(x\succeq^M_P y\succeq^M_P z\). If \(x=y\) or \(y=z\),
one premise is already the desired comparison. Otherwise both premises
are strict, and the strict transitivity just proved yields
\(x\succ^M_P z\), hence its weak version. Thus the output belongs
to \(\mathcal W(X)\), and distinct points are never socially
indifferent. This proves the stated map on every admissible profile.
The proof used only the axis order and finitely many voters; it applies
in particular to all of \([0,1]\), not merely a finite sample from it.
\end{proof}
\end{annexitem}

\begin{annexitem}{annex:p17.median-peak-winner}{thm:p17-median-peak-winner}{The two majority coalitions around the middle peak}
\begin{proof}
Sort the finite list of peaks, with repetitions, and write
\(m=p_{(k+1)}\). At least \(k+1\) peaks are at least \(m\),
since the entries from positions \(k+1\) through \(2k+1\)
have that property. At least \(k+1\) peaks are at most \(m\),
by the entries in positions one through \(k+1\).

If \(x<m\), every voter with \(p_i\geq m\) satisfies
\(p_i\geq m>x\). The first branch of single-peakedness gives
\(m\succ_i x\). There are at least \(k+1\) such voters, so
\(m\succ^M_P x\). If \(x>m\), every voter with
\(p_i\leq m\) satisfies \(x>m\geq p_i\), and the second
branch gives \(m\succ_i x\), again for at least \(k+1\)
voters. These cover every alternative distinct from \(m\), proving
that \(m\) is a strict Condorcet winner independently of the
preceding proof of transitivity.

Two distinct strict Condorcet winners would each strictly defeat the
other, contradicting the asymmetric pair-count comparison. Thus this
winner is unique. It is also a weak winner. Every other alternative
loses strictly to \(m\), so fails the definition of a weak winner.
This proves uniqueness in that sense too. In the already established
majority ordering, defeating every other alternative is exactly the
assertion that \(m\) is greatest.

To prove single-peakedness of the entire social ordering, take
\(m\geq b>c\). Every voter with \(p_i\geq m\) has
\(p_i\geq b>c\), hence \(b\succ_i c\); there are at least
\(k+1\) such voters. Thus \(b\succ^M_P c\). If
\(c>b\geq m\), use instead the at least \(k+1\) voters with
\(p_i\leq m\), for whom \(c>b\geq p_i\) gives the same
strict preference for \(b\). These are precisely the two defining
branches for a single-peaked social ordering with peak \(m\).

Repetition of a peak can only increase the sizes of these weak-inequality
coalitions, so it causes no exception. At an endpoint of \([0,1]\),
one set of comparisons is empty and the other argument still applies.
When \(k=0\), the only peak is the median and its single voter
supplies every required strict comparison. On a singleton \(X\)
there is no rival, and the stated winner and ordering conclusions
hold with the given definitions.
\end{proof}
\end{annexitem}

\begin{annexitem}{annex:p17.even-median-interval}{thm:p17-even-median-interval}{All weak winners and the coincident-central-peaks case}
\begin{proof}
Let \(l=p_{(k)}\), \(u=p_{(k+1)}\), and first take
\(w\in[l,u]_X\). If \(x<w\), at least \(k\) voters have
peaks at least \(u\), and hence at least \(w\). Each strictly
prefers \(w\) to \(x\). Therefore
\(V_P(w,x)\geq k\) and \(V_P(x,w)\leq k\), which is
\(w\succeq^M_P x\). If \(x>w\), the at least \(k\) voters
with peaks at most \(l\) all strictly prefer \(w\) to \(x\),
giving the same weak majority conclusion. Reflexivity handles
\(x=w\). Thus every member of the interval is a weak winner.

If \(w<l\), the sorted peak positions \(k,k+1,\ldots,2k\)
give at least \(k+1\) voters with \(p_i\geq l>w\).
All strictly prefer \(l\) to \(w\), so \(l\succ^M_P w\)
and \(w\) is not a weak winner. If \(w>u\), the positions
\(1,\ldots,k+1\) give at least \(k+1\) voters with
\(p_i\leq u<w\); hence \(u\succ^M_P w\), and again
\(w\) is not a weak winner. These exhaust the alternatives outside
\([l,u]_X\), proving the exact characterization.

When \(l=u=p\), there are at least \(k+1\) voters with
peaks at least \(p\), and at least \(k+1\) with peaks at most
\(p\), because both central entries equal \(p\). The odd-case
coalition argument, now with this strict majority of \(2k\), proves
that \(p\) strictly defeats every distinct alternative. Conversely,
when \(l<u\), both \(l,u\) are distinct weak winners. A strict
winner would have to defeat every other weak winner, which is impossible:
a weak winner by definition has no strict defeat. Formally, if the
proposed strict winner equals \(l\), compare it with \(u\);
otherwise compare it with \(l\). Thus there is no strict winner.
None of these arguments asserts transitivity of the even-electorate
weak majority relation.
\end{proof}
\end{annexitem}

\begin{annexitem}{annex:p17.single-peaked-algorithm}{thm:p17-single-peaked-algorithm}{A complete finite certificate and the operation counts}
\begin{proof}
The supplied axis labels have distinct indices \(1,\ldots,q\).
For each input list, an array of occurrence counts checks its length,
that every index lies in this range, and that each occurs exactly once.
A missing, repeated or out-of-range index rejects the claimed encoding.
For a valid permutation, assigning \(r_i(j)\) to the position of
index \(j\) gives the strict ordering
\(a_j\succ_i a_t\) exactly when \(r_i(j)<r_i(t)\).
The rank values are distinct integers, so this relation is irreflexive,
transitive and compares every pair. Its greatest alternative is its
unique rank-one entry \(a_{h_i}\).

If that ordering is single-peaked, the structural lemma makes its peak
\(a_{h_i}\). For \(j<h_i\), the point \(a_{j+1}\) is strictly
closer toward that peak on its left side in the order sense:
\(a_j<a_{j+1}\leq a_{h_i}\). The definition gives
\(a_{j+1}\succ_i a_j\), or \(r_i(j)>r_i(j+1)\).
For \(h_i\leq j<q\), the other branch gives
\(a_j\succ_i a_{j+1}\), or \(r_i(j)<r_i(j+1)\).
Thus an admissible ordering passes every displayed check.

Conversely suppose all adjacent checks pass. For any
\(s<t\leq h_i\), chaining the left-side rank inequalities gives
\(r_i(s)>r_i(t)\), hence \(a_t\succ_i a_s\).
For any \(h_i\leq s<t\), chaining the right-side inequalities
gives \(r_i(s)<r_i(t)\), hence \(a_s\succ_i a_t\).
These are all comparisons required by the two single-peaked branches,
including comparisons with the peak. Therefore the ordering is
single-peaked. A failed left adjacent check reverses one required
comparison on the left, and a failed right check reverses one on the
right. Since every possible peak is already forced to be the rank-one
alternative, no different choice of peak can repair that witness.
For \(q=1\), there are no adjacent checks and the sole ordering
is correctly accepted. This proves necessity and sufficiency of the
finite certificate for each voter and hence for the profile.

The peak histogram counts the entries at each axis index. Scanning
indices from left to right and adding their counts locates the first
cumulative total reaching any specified sorted position. This is
exactly that position in the sorted peak list, including repeated
peaks. It therefore yields \(p_{(k+1)}\) for \(n=2k+1\) or
\(p_{(k)},p_{(k+1)}\) for \(n=2k\). The proved median theorems
give the stated winner outputs.

For the full pairwise matrix, each comparison of \(r_i(j),r_i(t)\)
adds one vote to exactly one direction. After all \(n\) comparisons
the counts are precisely the defining \(V_P(a_j,a_t)\) and its
complement to \(n\). When \(n\) is odd, the earlier majority
theorem makes this pairwise relation a strict ordering. If \(x\)
is above \(y\) in that ordering, every alternative defeated by
\(y\) is also defeated by \(x\), by transitivity. In addition,
\(x\) defeats \(y\), whereas \(y\) does not defeat itself
or \(x\). Thus the victory count of \(x\) is strictly larger.
There are \(q\) alternatives and each victory count is an integer
between zero and \(q-1\). Since every pair has distinct counts,
the counts are exactly those \(q\) integers. Placing alternatives
into an array indexed by their victory counts and reading it downward
therefore prints precisely the majority ordering.

In the stated random-access model, reading and validating \(n\)
lists of length \(q\), building their rank arrays, and performing
\(q-1\) adjacent comparisons per voter use \(O(nq)\) operations.
Initializing the \(q\)-cell peak histogram, recording its \(n\)
peaks and scanning it take \(O(n+q)\) operations once the peaks
are known. There are exactly \(q(q-1)/2\) unordered pairs, and
the matrix step makes \(n\) rank comparisons per pair. Thus the
claimed individual-comparison count is exact. Reading its \(O(q^2)\)
entries supplies all victory counts, and the final count-indexed array
uses \(O(q)\) further operations. Together these costs are
\(O(nq^2)\), with \(n,q\geq1\), including validation.
The indices and counts are bounded by \(q\) and \(n\), respectively;
the unit-cost statement assumes words sufficient to store these values.
It is not a bit-complexity claim for arbitrary real or infinite inputs.
\end{proof}
\end{annexitem}

\begin{annexitem}{annex:p17.single-peaked-boundaries}{thm:p17-single-peaked-boundaries}{The five-voter computation, a cycle and nontransitive social ties}
\begin{proof}
For the worked five-voter profile, voters 1 and 5 are strictly monotone
along the axis toward their endpoint peaks. Voter 2's left side is
\(b\succ_2 a\), and its right side is
\(b\succ_2 c\succ_2 d\succ_2 e\). Voter 3's left side is
\(c\succ_3 b\succ_3 a\), and its right side is
\(c\succ_3 d\succ_3 e\). Voter 4's left side is
\(d\succ_4 c\succ_4 b\succ_4 a\), and its right side is
\(d\succ_4 e\). These verify all same-side comparisons and
therefore the single-peaked conditions.

For each pair with its first member to the left on the axis, the
voters preferring that member are exactly as follows:
\[
 \begin{array}{c|c|c}
 (x,y)&\{i:x\succ_i y\}&V_P(x,y)\\\hline
 (a,b)&\{1\}&1\\
 (a,c)&\{1,2\}&2\\
 (a,d)&\{1,2\}&2\\
 (a,e)&\{1,2,4\}&3\\
 (b,c)&\{1,2\}&2\\
 (b,d)&\{1,2,3\}&3\\
 (b,e)&\{1,2,3,4\}&4\\
 (c,d)&\{1,2,3\}&3\\
 (c,e)&\{1,2,3,4\}&4\\
 (d,e)&\{1,2,3,4\}&4
 \end{array}
\]
The reverse counts are five minus the displayed counts, verifying
every entry of the matrix. Reading the majority direction of every
pair gives exactly \(c\succ^M b\succ^M d\succ^M a\succ^M e\).
The peak list gives median \(c\). The table has three votes for
\(a\) over \(e\), while voter 3's given ordering puts \(e\)
above \(a\). Thus the stated failure of agreement with that voter's
whole ordering is an actual pairwise disagreement.

In the three-voter counterexample, voters 1 and 3 prefer \(x\) to
\(y\), voters 1 and 2 prefer \(y\) to \(z\), and voters 2
and 3 prefer \(z\) to \(x\). Each direction consequently wins
by two to one, producing the claimed cycle. In voter 3's ordering,
both \(x\) and \(z\) are above the axis-middle alternative
\(y\), so the structural lemma excludes it from the single-peaked
domain. The example does not refute the restricted-domain theorem.

For the four-voter example, the first two voters have peak \(x\),
the third has peak \(y\), and the fourth has peak \(z\).
The endpoint rankings move strictly toward their endpoint peaks, and
the middle-peaked ranking \(y\succ z\succ x\) satisfies both
of its two same-side requirements. Thus all inputs are single-peaked.
The two \(x\)-peaked voters prefer \(x\) to both other alternatives,
while the other two voters put \(x\) below both. Hence the counts
are two to two for each comparison involving \(x\). Both
\(x\)-peaked voters and the \(y\)-peaked voter prefer \(y\)
to \(z\), so \(y\) defeats \(z\) by three votes to one.
These counts give \(z\succeq^M x\), \(x\succeq^M y\),
and failure of \(z\succeq^M y\), proving nontransitivity of
the complete weak relation. Alternatives \(x,y\) have no strict
defeat; \(z\) loses to \(y\). Thus the weak winners are exactly
\(x,y\). Neither defeats the other strictly, so there is no strict
winner. The sorted peaks are \(x,x,y,z\), giving the stated
central pair and agreement with the even-median theorem.

This parity failure also occurs on the source's whole interval.
For a specified \(p\in[0,1]\), define the real distance function
\[
 d_p(t)=\max\{t-p,p-t\}\qquad(t\in[0,1]).
\]
Order alternatives by smaller \(d_p(t)\), and when two distances agree,
prefer the larger value of \(t\). Formally, \(s\) is above \(t\) when
\[
 d_p(s)<d_p(t)\quad\hbox{or}\quad
       (d_p(s)=d_p(t)\ \hbox{and}\ s>t).
\]
This is a strict ordering: distinct points have either unequal distances
or a strict numerical tie-break. In a chain with \(s\) above \(t\)
and \(t\) above \(v\), the distances satisfy
\(d_p(s)\leq d_p(t)\leq d_p(v)\). If either inequality is strict,
then \(s\) is above \(v\) by distance. If both are equal, the
tie-break comparisons give \(s>t>v\), hence again \(s\) is above
\(v\). This proves transitivity, and the definition also excludes
self-comparisons. If \(t\leq p\), the definition gives \(d_p(t)=p-t\);
if \(t\geq p\), it gives \(d_p(t)=t-p\). On either side of \(p\),
movement strictly toward \(p\) therefore strictly decreases distance, so the source's
single-peaked condition holds. With two voters having \(p=0\),
one \(p=1/2\), and one \(p=1\), their restrictions to
\(x=0,y=1/2,z=1\) are exactly the four specified rankings.
The same three pair comparisons therefore witness failure of
transitivity on the full interval. The tie-break removes individual
indifferences; the calculation concerns the declared strict-input domain.
\end{proof}
\end{annexitem}

\begin{annexitem}{annex:p17.single-peaked-arrow}{thm:p17-single-peaked-arrow}{Restricted-domain aggregation and the missing unrestricted premise}
\begin{proof}
The odd-majority theorem has already established completeness and
transitivity on every admissible profile. If all voters strictly
prefer \(x\) to \(y\), their counts are \(n,0\); hence the
social comparison is strict in the same direction, proving Pareto.
Permuting voter identities permutes the contributors to every count
without changing its value. The output relation is consequently
unchanged, which is anonymity. If two admissible profiles give each
voter the same comparison of \(x,y\), their counts in each
direction coincide, so both social weak comparisons, including the
strict direction they specify, coincide. This proves IIA with its
two-profile quantifiers ranging over the restricted domain.

To refute dictatorship, fix any proposed voter \(j\). Give voter
\(j\) the strictly decreasing preference along \([0,1]\): it
prefers \(s\) to \(t\) exactly when \(s<t\), with peak zero.
Give every other voter the strictly increasing preference, which
prefers \(s\) to \(t\) exactly when \(s>t\), with peak one.
These are strict single-peaked orderings on the entire interval.
Voter \(j\) prefers zero to one, but the \(n-1\) others prefer
one to zero. Since \(n\geq3\), \(n-1>1\), so the strict
majority comparison contradicts \(j\)'s comparison. This provides
an admissible refuting profile for every fixed candidate \(j\).
Therefore no voter dictates all profiles and pairs.

The domain restriction is proper. For example, a strict ordering
on \([0,1]\) can put \(0\) first, \(1\) second and \(1/2\)
third, then put all remaining alternatives below these three, ordering
that remainder by \(s\succ t\) exactly when \(s>t\). Concatenating this finite
ordered block and that strictly ordered remainder gives a strict
ordering: comparisons within each block are transitive and every
cross-block comparison points from the first block to the second.
Its restriction to \(0<1/2<1\) ranks the middle point below both
endpoints. It therefore fails single-peakedness by the structural
lemma. A profile containing this ordering belongs to the unrestricted
strict-profile space but not to \(\mathcal D_{\mathrm{SP}}\).

Arrow's theorem assumes a function defined on every unrestricted
strict profile. The present function is defined on the proper subset
\(\mathcal D_{\mathrm{SP}}\), so that premise is absent. The
distinction also appears in the proof: placing an interior reference
alternative at the bottom below alternatives on both sides creates
the forbidden valley, so the unrestricted extreme-profile
constructions need not remain admissible here. The other properties
proved above do not fill the missing domain premise. There is
therefore no contradiction with the stated Arrow theorem. For one
voter, every pair's majority comparison is that voter's own comparison,
and the output is its reflexive closure, exactly as claimed.
\end{proof}
\end{annexitem}
